% Chapter 2: Planning and Learning for Economists
% A narrative introduction to RL theory for PhD economists
% Frames RL as asymptotic approximations of the Bellman fixed-point operator
% Primary references: Bertsekas (2022, 2024), Puterman (1994), Agarwal et al. (2021)

\section{The Theory of Reinforcement Learning}
\label{sec:planning_learning}

\subsection{The Geometry of Dynamic Programming}
\label{sec:newton_connection}

Value iteration (VI) and policy iteration (PI) are the workhorses of dynamic programming. VI applies the Bellman operator repeatedly until convergence; PI alternates between \emph{policy evaluation} (solving a linear system) and \emph{policy improvement} (taking the greedy action). PI converges faster. Why? The answer reveals a connection between dynamic programming and numerical optimization. Policy iteration is Newton's method applied to the Bellman equation.

\subsubsection{Value Iteration as Picard Iteration}

Consider the Bellman optimality operator $T$ acting on value functions:
\begin{equation}
(TV)(s) = \max_{a \in \mathcal{A}} \left\{ r(s,a) + \gamma \sum_{s' \in \mathcal{S}} P(s'|s,a) V(s') \right\}.
\label{eq:bellman_operator}
\end{equation}
This operator is nonlinear due to the $\max$. Value iteration applies $T$ repeatedly: $V_{k+1} = TV_k$. Since $T$ is a $\gamma$-contraction in the supremum norm \citep{denardo1967}, Banach's fixed-point theorem guarantees $\|V_k - V^*\|_\infty \leq \gamma^k \|V_0 - V^*\|_\infty$.\footnote{This is the Contraction Mapping Theorem. The identical mathematical structure governs convergence of value function iteration in consumption-savings models, competitive equilibrium computation, and Bellman equation solution.} This is Picard iteration, with linear convergence at rate $\gamma$.\footnote{Picard iteration is $x_{k+1} = f(x_k)$ for finding roots of $x = f(x)$. When $f$ is a contraction, convergence is geometric. The rate $\gamma$ means each iteration reduces the error by a fixed proportion; more patient agents (higher $\gamma$) face slower convergence because the operator contracts less per step.} The iteration count to reduce error by a factor of $\delta$ is $k = \log(\delta) / \log(1/\gamma)$ \citep{bertsekas1996}.

\subsubsection{Policy Iteration as Newton's Method}

Policy iteration takes a different approach. At the current value estimate $\tilde{V}$, define the greedy policy $\tilde{\pi}(s) = \argmax_a \{ r(s,a) + \gamma \sum_{s'} P(s'|s,a) \tilde{V}(s') \}$. The \emph{policy evaluation} step solves the linear fixed-point equation $V = T^{\tilde{\pi}} V$ exactly, where $T^{\tilde{\pi}}$ is the policy-specific Bellman operator:
\begin{equation}
(T^{\tilde{\pi}} V)(s) = r(s, \tilde{\pi}(s)) + \gamma \sum_{s'} P(s'|s,\tilde{\pi}(s)) V(s').
\end{equation}

The geometric structure, formalized by \citet{puterman1979}, is as follows: the linear operator $T^{\tilde{\pi}}$ is a supporting hyperplane to the nonlinear operator $T$ at the current iterate.\footnote{A supporting hyperplane to a convex function at $x_0$ is a linear function $\ell$ with $\ell(x_0) = f(x_0)$ and $\ell(x) \leq f(x)$ everywhere. In plainer terms: $T^{\tilde{\pi}}$ is the tangent-line approximation to $T$ from elementary calculus, extended to function spaces. The policy operator $T^{\tilde{\pi}}$ plays this role for the Bellman operator.} Specifically, the operators satisfy tangency: $T^{\tilde{\pi}} \tilde{V} = T\tilde{V}$, so the linearization agrees with the nonlinear operator at the current iterate. They also satisfy support: $T^{\tilde{\pi}} V \leq TV$ for all $V$, meaning the linear operator lies weakly below the nonlinear one everywhere, just as a tangent line lies below a convex function. Policy evaluation solves for the fixed point of this linearization exactly. This is precisely the structure of Newton's method; linearize the nonlinear equation at the current point, solve the linearized system, and iterate.\footnote{The Newton interpretation of policy iteration has precursors in \citet{kleinman1968} for Riccati equations in linear-quadratic control and \citet{pollatschek1969} for stochastic games.}\footnote{Algebraically: consider finding the root of $F(V) = V - TV = 0$. For a fixed policy $\pi$, the Bellman operator is affine: $T^\pi V = r^\pi + \gamma P^\pi V$. The Jacobian of the residual $F$ with respect to $V$ is $(I - \gamma P^{\tilde{\pi}})$, where $\tilde{\pi}$ is the greedy policy at the current iterate. Policy evaluation solves $(I - \gamma P^{\tilde{\pi}})V = r^{\tilde{\pi}}$, which is algebraically the Newton step $V_{k+1} = V_k - [\nabla F(V_k)]^{-1} F(V_k)$. Because the Bellman operator is affine within regions where the greedy policy does not change, Newton's method converges in finitely many steps for tabular MDPs.}

The consequence is quadratic convergence near the optimum. While VI requires $k = \log(100) / \log(1/\gamma)$ iterations to reduce error by a factor of 100 \citep{bertsekas1996}, PI typically converges in 5--10 iterations regardless of $\gamma$, an empirical observation consistent with the quadratic convergence rate established by \citet{puterman1979}.\footnote{\citet{ye2011} proves PI is strongly polynomial with iteration count $O\!\left(\frac{|\mathcal{S}||\mathcal{A}|}{1-\gamma}\log\frac{|\mathcal{S}|}{1-\gamma}\right)$, resolving a long-standing conjecture. This bound is for fixed $\gamma$; \citet{fearnley2010} constructs examples requiring exponentially many iterations when $\gamma$ is allowed to vary with $|\mathcal{S}|$.}\footnote{Convergence is superlinear in general. The ``quadratic'' rate requires regularity conditions: \citet{santos2004} show that for continuous-state problems (the norm in economics), the convergence order is approximately 1.5 unless the value function is piecewise linear and concave, in which case true quadratic convergence obtains.} At $\gamma = 0.90$, VI needs 44 iterations while PI needs only 5--8; at $\gamma = 0.95$, VI requires 90 versus 5--8; at $\gamma = 0.99$, VI needs 459 iterations while PI still converges in 5--10.

\citet{bertsekas2022newton} extends this interpretation to a broad class of dynamic programming problems. The Newton structure applies whenever the Bellman operator can be written as a pointwise maximum over linear operators: $T = \max_\pi T^\pi$. This includes not only infinite horizon problems with discounting but also optimal stopping problems (job search, option exercise), average cost optimization (inventory, queueing), and minimax formulations for adversarial settings.\footnote{These problems appear under different names such as ``stochastic shortest path'' for optimal stopping, ``average-cost MDP'' for long-run average optimization, and ``model predictive control'' (or receding-horizon control) for finite-horizon replanning.} The practical implication is that algorithms with policy-improvement structure (evaluate a policy exactly or approximately, then improve) inherit Newton-like convergence behavior, while pure value-iteration methods (apply $T$ directly) are limited to linear convergence.\footnote{\citet{blackwell1965} proves that for discounted MDPs with finite state and action spaces, a stationary deterministic policy $\pi: \mathcal{S} \to \mathcal{A}$ exists that is optimal for all initial states simultaneously. This ``uniform optimality'' has three implications: (1) the search space reduces from history-dependent or stochastic policies to static maps, justifying neural networks that condition only on current state; (2) the optimal policy is independent of the initial distribution $d_0$, so changing where episodes start does not require retraining; (3) PI and VI are guaranteed to converge to the same globally optimal policy regardless of initialization.}

\subsection{Value Learning Methods}
\label{sec:stochastic_approx}

\subsubsection{Stochastic Approximation Foundations}

When $P$ is unknown, a single sampled transition $(s, a, r, s')$ can replace the expectation $\mathbb{E}_{s' \sim P(\cdot|s,a)}[V(s')]$. The mathematical foundation is stochastic approximation, developed by \citet{Robbins1952}.\footnote{The modern theory of stochastic approximation, including convergence rates and the ODE (ordinary differential equation) method for analyzing iterates, is developed in \citet{kushner1978} and \citet{borkar2000}. The ODE method shows that the expected trajectory of the stochastic iterates tracks the solution of a deterministic differential equation $\dot{x} = -g(x)$, providing stability conditions via Lyapunov theory.} Consider the problem of finding $x^*$ such that $g(x^*) = 0$, where $g$ cannot be evaluated directly but one can observe noisy samples $g(x) + \epsilon$. The Robbins-Monro iteration is:
\begin{equation}
x_{t+1} = x_t - \alpha_t [g(x_t) + \epsilon_t],
\label{eq:robbins_monro}
\end{equation}
where $\epsilon_t$ is zero-mean noise. Under two conditions on the step sizes, this converges to $x^*$ with probability one. The conditions are $\sum_{t=0}^\infty \alpha_t = \infty$ (sufficient exploration, ensuring that learning never ceases) and $\sum_{t=0}^\infty \alpha_t^2 < \infty$ (diminishing noise, ensuring the variance of cumulative updates is finite). The canonical choice $\alpha_t = 1/(t+1)$ satisfies both conditions.

\subsubsection{Q-Learning and SARSA}

Q-learning \citep{WatkinsDayan1992} is Robbins-Monro applied to the Bellman equation for action-value functions.\footnote{The ``Q'' in Q-learning stands for ``quality,'' following \citet{WatkinsDayan1992}, who used $Q(s,a)$ to denote the quality (expected return) of taking action $a$ in state $s$. The term ``Q-factor'' is used interchangeably with ``action-value function'' throughout the RL literature.} Define the Q-factor Bellman operator:
\begin{equation}
(FQ)(s,a) = r(s,a) + \gamma \sum_{s'} P(s'|s,a) \max_{a'} Q(s',a').
\end{equation}
The Q-learning update, upon observing transition $(s_t, a_t, r_t, s_{t+1})$, is:
\begin{equation}
Q(s_t, a_t) \leftarrow Q(s_t, a_t) + \alpha_t \left[ r_t + \gamma \max_{a'} Q(s_{t+1}, a') - Q(s_t, a_t) \right].
\label{eq:qlearning}
\end{equation}

The logic of Q-learning is best understood as a Monte Carlo approximation of the Bellman contraction. The true Bellman operator involves an integral over the transition distribution, $(FQ)(s,a) = r + \gamma \int \max_{a'} Q(s',a') \, dP(s'|s,a)$. Since $P$ is unknown, this integral cannot be computed analytically. However, a sample transition $(s,a,r,s')$ acts as a single-point Monte Carlo estimate of this integral. The Q-learning update is simply an exponential moving average (with weight $\alpha_t$) between the current estimate and this noisy Monte Carlo target. Because $F$ is a $\gamma$-contraction in the supremum norm, the expected update drives the estimate toward the fixed point $Q^*$, provided the noise in the Monte Carlo sample averages out over time (which the Robbins-Monro conditions ensure) \citep{tsitsiklis1994}.\footnote{Q-learning can be viewed as root-finding for the expected Bellman residual: the goal is to find parameters such that $\mathbb{E}_{(s,a,r,s')}[\delta_t] = 0$, where $\delta_t = r + \gamma \max_{a'} Q(s',a') - Q(s,a)$ is the temporal difference error. The update is a stochastic gradient step on the mean-squared Bellman error, but with a critical distinction: the gradient is ``semi-gradient'' because the target $\max_{a'} Q(s',a')$ is treated as a fixed constant rather than a function of the parameters being updated. This simplifies computation but disconnects the update from true gradient descent, requiring the specific stability conditions of \citet{tsitsiklis1994}.}

Convergence requires two conditions. Exploration (visiting all state-action pairs infinitely often) ensures identification; every state-action pair must be sampled to learn its value. The Robbins-Monro step-size conditions ($\sum_t \alpha_t = \infty$, $\sum_t \alpha_t^2 < \infty$) balance tracking versus noise suppression. \citet{WatkinsDayan1992} and \citet{jaakkola1994} formalize these; \citet{tsitsiklis1994} provides the general framework.

The choice of step-size schedule has quantitative consequences: \citet{evendar2003} show that polynomial schedules $\alpha_t = 1/t^\omega$ with $\omega \in (1/2, 1)$ achieve convergence rate $O(1/t^{1-\omega})$, creating an explicit tradeoff between speed and stability. Recent work by \citet{li2024minimax} establishes that Q-learning with variance-reduced updates achieves minimax-optimal sample complexity $\tilde{O}(|\mathcal{S}||\mathcal{A}|/(1-\gamma)^3\epsilon^2)$, matching information-theoretic lower bounds.\footnote{Vanilla Q-learning (without variance reduction) has tight complexity $\tilde{\Theta}(|\mathcal{S}||\mathcal{A}|/(1-\gamma)^4\epsilon^2)$, worse by a factor of $1/(1-\gamma)$ due to maximization bias inflating variance. The optimal cubic rate requires variance-reduced updates \citep{wainwright2019,sidford2018}. By contrast, naive model-based RL (estimate $\hat{P}$ from samples, then solve by planning) achieves the optimal $(1-\gamma)^{-3}$ rate with no special tricks \citep{agarwal2020}, illustrating the statistical cost of discarding transition structure.} Model-free learning is possible. The optimal value function can be found without ever estimating the transition probabilities.\footnote{Model-free methods are essential when (a) the environment is a physical system with dynamics too complex to write down, or (b) the agent learns directly from interaction. Note that ``model-free'' does not mean ``atheoretical.'' The agent does not store $P(s'|s,a)$ explicitly, but the Q-function serves as an implicit model encoding long-run consequences.}

SARSA provides an on-policy variant.\footnote{SARSA is named for the quintuple $(S_t, A_t, R_t, S_{t+1}, A_{t+1})$ used in each update \citep{rummery1994}. Convergence requires the GLIE (Greedy in the Limit with Infinite Exploration) condition: the behavior policy must explore all actions infinitely often while converging to a greedy policy. $\varepsilon$-greedy with $\varepsilon_t \to 0$ satisfies this.} Instead of taking the maximum over next actions, SARSA uses the action actually taken:
\begin{equation}
Q(s_t, a_t) \leftarrow Q(s_t, a_t) + \alpha_t \left[ r_t + \gamma Q(s_{t+1}, a_{t+1}) - Q(s_t, a_t) \right].
\end{equation}
This solves for the value function of the behavior policy $\pi$ rather than the optimal policy. \citet{singh2000} prove convergence under the same step-size conditions, provided the behavior policy converges to a stationary distribution. Q-learning is noisy value iteration on Q-factors; SARSA is noisy \emph{policy evaluation}.\footnote{\citet{bhandari2021finite} provide finite-time analysis of TD learning, showing that the convergence rate depends on the mixing time of the Markov chain under the behavior policy. Faster mixing (less serial correlation in the state sequence) yields faster convergence.} The Robbins-Monro conditions ensure that the noise averages out faster than the signal decays, allowing asymptotic convergence despite using only single-sample estimates.

\subsubsection{AlphaZero and the Bellman Contraction}

Recall the AlphaGo Zero system from Section~\ref{subsubsec:alphago_zero}, where a single neural network $f_\theta(s) = (\mathbf{p}, v)$ outputs a prior policy $\mathbf{p}$ and a value estimate $v$ for any board position $s$. During play, the network does not act alone. Monte Carlo Tree Search runs 1,600 simulated games from the current position, using $v$ to evaluate leaf nodes and $\mathbf{p}$ to guide which branches to explore.\footnote{A \emph{rollout} simulates one or more trajectories forward from the current state to estimate the value of taking a particular action. The base policy $\mu$ provides a default behavior; the rollout policy improves on $\mu$ by computing one step of lookahead. In practice, from state $s$, try each action $a$, simulate forward using $\mu$ for $H$ steps, average the returns, and pick the action with highest estimated value. MCTS is a sophisticated rollout that grows a search tree rather than simulating independent trajectories. Rollout requires a simulator (generative model) that can be queried from arbitrary states, a stronger assumption than the trajectory-based access of Q-learning and SARSA.} This is precisely the structure from Section~\ref{sec:newton_connection}. The network provides an approximate value function $\tilde{V}$, and MCTS applies the Bellman operator $T$ through lookahead, executing one step of policy improvement.\footnote{Bertsekas uses cost-minimization notation throughout his work, writing $\min$ where standard RL uses $\max$ and defining value as accumulated cost rather than reward. I translate to the reward-maximization convention used elsewhere in this chapter; the mathematics are equivalent with reversed inequalities.} \citet{bertsekas2021lessons} proves the cost improvement property for this construction: $V^{\tilde{\pi}}(s) \geq V^{\mu}(s)$ for all $s$, with strict inequality unless the base policy $\mu$ is already optimal. Each move of AlphaZero \citep{Silver2018} is a Newton step in the sense of Section~\ref{sec:newton_connection}; the network supplies the current iterate, and MCTS solves the linearized subproblem through deep search. Table~\ref{tab:pi_alphazero} makes the correspondence explicit.

\begin{table}[h]
\centering
\small
\begin{tabular}{lll}
\hline
Step & Policy Iteration & AlphaZero \\
\hline
Initialize & Arbitrary $\pi_0$ & Random network $f_\theta$ \\
Evaluate & Solve $V = T^{\pi_k} V$ exactly & Network value head $v \approx V^{\pi_k}$ \\
Improve & $\pi_{k+1} = \argmax_a \{r + \gamma P V^{\pi_k}\}$ & MCTS simulations from $v$ \\
Iterate & Repeat until $\pi$ is stationary & Retrain $f_\theta$ on self-play outcomes \\
\hline
\end{tabular}
\caption{Policy iteration and AlphaZero follow the same evaluate-improve loop. The network provides approximate policy evaluation; MCTS provides exact (within the search tree) policy improvement.}
\label{tab:pi_alphazero}
\end{table}

Since the Bellman operator $T$ is a $\gamma$-contraction (Section~\ref{sec:newton_connection}), each step of lookahead shrinks the terminal approximation error by a factor of $\gamma$. An $H$-step lookahead with terminal approximation $\tilde{V}$ satisfies $\|V^{\tilde{\pi}} - V^*\|_\infty \leq \frac{2\gamma^H}{1-\gamma} \|\tilde{V} - V^*\|_\infty$.\footnote{An $H$-step lookahead builds a tree of depth $H$, evaluating all possible action sequences up to that horizon, with the terminal values $\tilde{V}(s')$ approximating the infinite future beyond the tree. At $H = 50$ and $\gamma = 0.99$, the error multiplier $2\gamma^H / (1-\gamma)$ is approximately 121, meaning moderate lookahead does not compensate for a poor base estimate. At $H = 500$ the multiplier drops to approximately 1.3, and sufficiently deep lookahead can absorb substantial approximation error. With finite state and action spaces, sufficiently deep lookahead from any approximation yields the optimal policy exactly \citep[Prop.~2.3.1]{bertsekas2022}.} \citet{bertsekas2021lessons} recognizes this as a Newton step on $V = TV$, inheriting the superlinear convergence from the policy-iteration-as-Newton framework. This explains why AlphaZero's online play (network plus 800 MCTS simulations per move) far exceeds the raw neural network: the network provides a rough starting point, and MCTS executes a high-quality contraction step through deep lookahead.\footnote{AlphaGo Zero uses 1,600 MCTS simulations per move for Go; the generalized AlphaZero algorithm uses 800 simulations per move across chess, shogi, and Go \citep[Table~S3]{Silver2018}.} The gap between network-only and network-plus-search is the difference between evaluating an approximate value function and applying the Bellman operator to it.

\subsection{The Central Challenge: The Deadly Triad}
\label{sec:deadly_triad}

The preceding results assume tabular representations, namely, a separate value for each state-action pair. In practice, state spaces are too large for lookup tables (Go has $10^{170}$ states; most economic models have continuous state variables). Practitioners must therefore combine three ingredients: function approximation (to handle large state spaces), bootstrapping (to learn from single transitions rather than complete episodes), and off-policy learning (to learn about the optimal policy while exploring, or to reuse old data). Each is desirable in isolation. Their interaction, known as the \emph{deadly triad} \citep[Ch.~11]{sutton2018}, is the central open problem in reinforcement learning theory.

\subsubsection{The Projected Bellman Operator}

With function approximation $V(s) \approx \phi(s)^\top \theta$, the parameter vector $\theta$ is shared across states. Updating $\theta$ to improve the value estimate at one state simultaneously changes the estimate at every other state. The algorithm can no longer apply the Bellman operator $T^\pi$ to each state independently; instead, it applies $T^\pi$ to compute a target, then \emph{projects} the result back onto the function space (the span of the features $\phi$). The composed operator is $\Pi T^\pi$, where $\Pi$ denotes this projection \citep{tsitsiklis1997}. Convergence of the approximate iteration $\theta_{k+1} = \Pi T^\pi \theta_k$ requires $\Pi T^\pi$ to be a contraction.

In the on-policy setting, $\Pi$ minimizes squared error weighted by $d^\pi$, the stationary distribution of the policy being evaluated: $\Pi V = \arg\min_{\hat{V} \in \text{span}(\Phi)} \|V - \hat{V}\|_{d^\pi}$, where $\|V\|_{d^\pi}^2 = \sum_s d^\pi(s) V(s)^2$. The Bellman operator $T^\pi$ is a $\gamma$-contraction in the same $d^\pi$-norm. Because both operators use the same norm, the projection $\Pi$ is \emph{orthogonal}: the residual $V - \Pi V$ is perpendicular to the approximation subspace. The Pythagorean theorem then gives $\|\Pi V\|_{d^\pi}^2 + \|V - \Pi V\|_{d^\pi}^2 = \|V\|_{d^\pi}^2$, so $\|\Pi V\|_{d^\pi} \leq \|V\|_{d^\pi}$.\footnote{The same argument in $\mathbb{R}^n$: projecting a vector onto a subspace never makes it longer. This is the geometric content of the Cauchy-Schwarz inequality. In the function-approximation setting, the ``vector'' is a value function, the ``subspace'' is the span of features, and ``length'' is the $d^\pi$-weighted $L^2$ norm.} The projection cannot expand distances. The composition therefore contracts:
\begin{equation}
\|\Pi T^\pi V_1 - \Pi T^\pi V_2\|_{d^\pi} \leq \underbrace{\|\Pi\|}_{\leq\, 1} \cdot \underbrace{\|T^\pi V_1 - T^\pi V_2\|_{d^\pi}}_{\leq\, \gamma \|V_1 - V_2\|_{d^\pi}} < \|V_1 - V_2\|_{d^\pi}.
\label{eq:on_policy_contraction}
\end{equation}
A unique fixed point $\Phi \theta^*$ exists, and TD($\lambda$) converges to it with probability one. The resulting approximation error satisfies $\|\Phi \theta^* - V^\pi\|_{d^\pi} \leq \frac{1-\lambda\gamma}{\sqrt{1-\gamma^2}} \|\Pi V^\pi - V^\pi\|_{d^\pi}$, bounding the TD solution's error by a multiple of the best possible approximation error \citep[Theorem~1]{tsitsiklis1997}.

\subsubsection{Why Off-Policy Learning Diverges}

In the off-policy setting, samples come from a behavior distribution $\mu \neq d^\pi$. The projection now minimizes error under $\mu$, but the Bellman operator $T^\pi$ still contracts in the $d^\pi$-norm. The two operators measure distance in different norms. The projection is no longer orthogonal in the $d^\pi$-norm; it is \emph{oblique}. Unlike orthogonal projections, oblique projections can expand distances: $\|\Pi_\mu\|_{d^\pi}$ can exceed 1. If $\|\Pi_\mu\|_{d^\pi} > 1/\gamma$, the expansion from projection overwhelms the $\gamma$-contraction from the Bellman operator, and the fixed-point iteration diverges.

This divergence is not overfitting. Overfitting occurs when the approximator memorizes training data at the expense of generalization; collecting more data helps. Divergence means the parameter vector $\theta$ grows without bound, producing arbitrarily large value estimates that bear no relation to the true values. More data does not help; the algorithm itself is unstable. The distinction matters because the remedies are entirely different: regularization and early stopping address overfitting, while the deadly triad requires structural changes to the algorithm.

\citet{Baird1995} constructed a six-state star MDP that makes this failure concrete. All rewards are zero, so the true value is $V^*(s) = 0$ for every state. A lookup table learns this immediately. The MDP has a star topology: states 1 through 5 each transition to state 6, and state 6 transitions to itself. Linear function approximation uses a shared weight $w_1$ across all states plus a state-specific weight, with $V(s) = 2w_1 + w_s$ for $s \in \{1,\ldots,5\}$ and $V(6) = 2w_1 - w_6$. Training samples all transitions equally often (uniform distribution, not $d^\pi$). The dynamics are as follows.\footnote{\citet{Baird1995} also presents an MDP variant with two actions per state, demonstrating that Q-learning diverges under the same mechanism. The key insight is identical: shared parameters create cross-state coupling that uniform sampling cannot counterbalance.} When $V(6)$ is large and positive, the TD target $\gamma V(6)$ exceeds $V(s)$ for states 1 through 5, producing positive TD errors that push $w_1$ upward. At state 6, the TD target is $\gamma V(6) < V(6)$, producing a negative TD error that pushes $w_1$ downward. But states 1 through 5 are each visited as often as state 6, so $w_1$ receives five upward pushes for every one downward push. The shared weight diverges to $+\infty$. The on-policy distribution would concentrate mass on state 6 (the absorbing state), counterbalancing the upward pressure; uniform sampling destroys this balance.\footnote{The gradient of the mean-squared Bellman error $\|Q - TQ\|^2$ would correct for this imbalance, but computing an unbiased gradient requires two independent next-state samples from the same $(s,a)$, since $\nabla\mathbb{E}[(Q - \mathbb{E}[r + \gamma V(s')])^2]$ involves the product $\mathbb{E}[\cdot] \cdot \nabla\mathbb{E}[\cdot]$. A single sample yields a biased estimate ($\mathbb{E}[XY] \neq \mathbb{E}[X]\mathbb{E}[Y]$). This ``double-sampling'' requirement is impractical in most environments \citep{Baird1995}. Practitioners instead use semi-gradient TD, which treats the bootstrap target as a constant and finds the projected fixed point $\Pi T^\pi \theta^*$ rather than the true Bellman fixed point. It is this semi-gradient structure that makes off-policy TD inherently vulnerable to the projection mismatch.}

Each element of the triad is individually necessary for divergence. Without function approximation (tabular), the projection is the identity and Q-learning's contraction applies directly. Without bootstrapping (Monte Carlo returns), targets are independent of current value estimates and the problem reduces to supervised regression. Without off-policy learning, samples come from $d^\pi$, the projection is orthogonal, and the Tsitsiklis-Van Roy convergence guarantee holds.

\subsubsection{Resolutions}
\label{sec:deadly_triad_resolutions}

Three classes of algorithms restore convergence, each neutralizing a different component of the triad.

\emph{Target networks} weaken bootstrapping. Instead of updating toward $r + \gamma Q(s'; \theta)$, where the target moves with each parameter update, DQN \citep{mnih2015} updates toward $r + \gamma Q(s'; \theta^-)$, where $\theta^-$ is a slowly-updated copy of the parameters.\footnote{Experience replay \citep{Lin1992} complements target networks by breaking temporal correlation in the training data. The two mechanisms address different sources of instability: target networks stabilize the bootstrap target, while replay stabilizes the sampling distribution.} The regression target becomes quasi-static, converting the coupled fixed-point problem into a sequence of supervised learning problems. \citet{zhang2021target} prove that this two-timescale scheme converges to a regularized TD fixed point with linear function approximation. \citet{Fellows2023} show that target networks recondition the Jacobian of the TD update: the spectral radius of the composed update operator depends on the target network update frequency $k$, and for sufficiently large $k$ the spectral radius drops below 1 even in off-policy settings with nonlinear function approximation.

\emph{Gradient TD methods} fix the projection mismatch. \citet{sutton2009gtd} reformulate the projected Bellman error as a saddle-point problem $\min_\theta \max_y L(\theta, y)$, yielding algorithms (GTD, GTD2, TDC) that perform true stochastic gradient descent on the mean-squared projected Bellman error.\footnote{The saddle-point formulation introduces auxiliary variables $y$ of the same dimension as $\theta$, doubling the parameter count and requiring a second learning rate. \citet{bhandari2021finite} provide finite-time convergence rates for these two-timescale algorithms.} These methods converge off-policy with linear function approximation because they eliminate the semi-gradient approximation that causes the norm mismatch.

\emph{Regularization} shrinks the projection operator. \citet{lim2024regularized} add an $\ell_2$ penalty $-\eta\theta$ to the Q-learning update. This changes the projection from $\Pi = X(X^\top D X)^{-1} X^\top D$ to $\Pi_\eta = X(X^\top D X + \eta I)^{-1} X^\top D$. As the regularization strength $\eta$ increases, the projection ``shrinks'' toward the origin. For sufficiently large $\eta$, $\gamma\|\Pi_\eta\| < 1$, restoring the contraction property. The algorithm converges to a biased but stable fixed point, with the bias controlled by $\eta$.

%\subsection{When Value-Based RL Works}
%\label{sec:when_value_works}

%The deadly triad identifies conditions for divergence. The natural question is under what conditions does value-based RL with function approximation succeed? Theoretical advances identify structural assumptions that restore convergence and yield finite-sample guarantees.

%\subsubsection{Concentrability and Data Coverage}

%The central result from finite-time analysis of Fitted Value Iteration \citep{MunosSzepesvari2008} is that statistical efficiency requires the sampling distribution $\mu$ to provide adequate coverage of the state-action space visited by near-optimal policies. The concept originates with \citet{munos2003} and was refined by \citet{farahmand2010}, who provide tighter bounds under smoothness assumptions. This is formalized through the concentrability coefficient $C$, which measures the worst-case ratio between the discounted state-action occupancy of any policy and the sampling distribution:
%\begin{equation}
%C = \max_{\pi} \sup_{s,a} \frac{d^\pi(s,a)}{\mu(s,a)}.
%\end{equation}
%When $C$ is bounded (e.g., $\varepsilon$-greedy behavior ensures $\mu(s,a) \geq \varepsilon/|\mathcal{A}|$ for all visited states), the bound is finite but may be exponentially large in the horizon. When $C$ is finite, FVI achieves error bounds scaling polynomially in $C$, the horizon $H = 1/(1-\gamma)$, and the approximation capacity of the function class. The concentrability condition is strong: it requires the offline dataset to cover all states the optimal policy might visit. This is the statistical cost of breaking the contraction principle---the data must compensate for the lost geometric convergence.

%\subsubsection{Structural Assumptions: Linear MDPs and Bellman Rank}

%The Linear MDP assumption \citep{jin2020} provides an alternative path by imposing structure on the dynamics rather than the data. It posits that transition probabilities and rewards are linear in a known feature map $\phi(s,a) \in \mathbb{R}^d$:
%\begin{equation}
%P(s'|s,a) = \langle \phi(s,a), \mu(s') \rangle, \quad r(s,a) = \langle \phi(s,a), \theta^* \rangle,
%\end{equation}
%where $\mu: \mathcal{S} \to \mathbb{R}^d$ is an unknown measure and $\theta^* \in \mathbb{R}^d$ is an unknown parameter vector. Under this assumption, the Bellman operator preserves linearity: applying $T$ to a linear Q-function yields a function that remains linear in the features. This closure property sidesteps the projection error that drives the deadly triad. \citet{jin2020} prove that LSVI-UCB achieves regret $\tilde{O}(\sqrt{d^3 H^3 T})$, polynomial in the feature dimension $d$ and horizon $H$, but independent of $|\mathcal{S}|$ and $|\mathcal{A}|$.

%The concept generalizes through Bellman rank \citep{Jiang2017}, which measures the complexity of a value function class under Bellman backups, and Eluder dimension \citep{wang2020eluder}, which captures how quickly exploration reduces uncertainty. If the Bellman rank is $M$, then sample complexity scales polynomially in $M$ rather than $|\mathcal{S}|$. Linear MDPs have Bellman rank at most $d$; other structured classes (block MDPs, factored MDPs, linear quadratic regulators)\footnote{Block MDPs group states into latent clusters with shared dynamics; factored MDPs decompose the state into independent components; LQR (linear-quadratic regulator) problems have linear dynamics and quadratic costs. All have low Bellman rank. The hierarchy from tabular through linear to general function approximation mirrors the progression from nonparametric to semiparametric to parametric models in econometrics: stronger structural assumptions yield sharper statistical guarantees.} also have low Bellman rank. These results delineate when function approximation is provably safe: polynomial sample complexity is achievable when the environment's dynamics lie within a low-complexity function class. Without such structure, \citet{du2021} and \citet{zanette2020} prove exponential lower bounds, showing that general function approximation is statistically intractable.

%The contrast with policy gradient methods is instructive. Policy gradients achieve global convergence without concentrability assumptions \citep[Theorem~5.3]{agarwal2021theory}, but require on-policy data and suffer high variance. Value-based methods can use off-policy data efficiently, but only when the data distribution satisfies coverage conditions or the MDP has low structural complexity. The choice between methods trades off data requirements against computational and statistical efficiency.


\subsection{Policy Learning Methods}
\label{sec:policy_gradient}

Value-based methods find fixed points of the Bellman operator. Policy-based methods parameterize the policy directly as $\pi_\theta(a|s)$ and maximize expected return $J(\theta) = \mathbb{E}_{\pi_\theta}[\sum_{t=0}^\infty \gamma^t R_t]$ by gradient ascent. This formulation sidesteps the Bellman equation entirely and frames reinforcement learning as constrained optimization.

\subsubsection{The Policy Gradient Theorem}

The policy gradient theorem, proved independently by \citet{williams1992} for the episodic case and \citet{SuttonMcAllester2000} for the general discounted setting, provides a tractable expression for the gradient:
\begin{equation}
\nabla_\theta J(\theta) = \mathbb{E}_{s \sim d^{\pi_\theta}, a \sim \pi_\theta} \left[ \nabla_\theta \log \pi_\theta(a|s) \, Q^{\pi_\theta}(s,a) \right],
\label{eq:policy_gradient}
\end{equation}
where $d^{\pi_\theta}(s) = (1-\gamma) \sum_{t=0}^\infty \gamma^t \mathbb{P}(s_t = s | \pi_\theta)$ is the discounted state visitation distribution. The fundamental econometric challenge in optimizing $J(\theta)$ is that this distribution depends on $\theta$ through the environment's dynamics. A naive derivative would require $\nabla_\theta d^{\pi_\theta}(s)$, which implies differentiating the unknown transition matrix $P(s'|s,a)$.

The policy gradient theorem sidesteps this entirely via a likelihood ratio (or score function) trick.\footnote{The score function $\nabla_\theta \log \pi_\theta(a|s)$ is the same mathematical object that appears in the Cram\'er-Rao bound and the score test in maximum likelihood estimation. The policy gradient is a covariance: $\nabla J = \text{Cov}_{d^{\pi_\theta} \times \pi_\theta}(\nabla \log \pi_\theta, Q^{\pi_\theta})$, measuring how sensitive the log-likelihood of the policy is to parameter changes, weighted by action quality.} The theorem transforms a sensitivity analysis problem (how does the system evolve?) into a simpler expectation problem (what is the correlation between the score $\nabla \log \pi$ and the value $Q$?). The gradient can be written as an expectation under the current policy, weighted by action-values, without requiring $\nabla_\theta d^{\pi_\theta}$. The transition dynamics $P(s'|s,a)$ do not appear; the gradient is estimable via sample averages from trajectories alone.

\subsubsection{REINFORCE and Variance Reduction}

REINFORCE \citep{williams1992} is the simplest policy gradient algorithm. Sample a trajectory $(s_0, a_0, r_0, s_1, \ldots)$, compute the return $G_t = \sum_{k=0}^\infty \gamma^k r_{t+k}$ from each time step, and update:
\begin{equation}
\theta \leftarrow \theta + \alpha \sum_t \nabla_\theta \log \pi_\theta(a_t|s_t) \, G_t.
\end{equation}
This is an unbiased estimator of $\nabla_\theta J(\theta)$, but its variance is high because a single trajectory provides a noisy estimate of $Q^{\pi_\theta}$. Despite high variance, REINFORCE converges to a globally optimal policy in the tabular setting.\footnote{The baseline $b(s)$ subtracted from $G_t$ reduces variance while preserving unbiasedness, since $\mathbb{E}[\nabla_\theta \log \pi_\theta(a|s) b(s)] = 0$ for any baseline independent of $a$.}

\subsubsection{Natural Policy Gradient and Gradient Domination}

Standard gradient descent treats all parameter directions equally. But small changes in $\theta$ can cause large changes in the policy distribution $\pi_\theta$. The natural policy gradient \citep{Kakade2001}, building on the natural gradient framework of \citet{amari1998natural}, accounts for this curvature by preconditioning with the Fisher information matrix.\footnote{The Fisher information $F(\theta) = \mathbb{E}[\nabla \log p \nabla \log p^\top]$ measures curvature of the log-likelihood and appears in the Cram\'er-Rao bound. Here it measures how policy distributions change with parameters.} The relationship between NPG and standard PG parallels that between GLS and OLS in econometrics: both account for the covariance structure of the estimating equations, yielding efficiency gains when the covariance is non-spherical.
\begin{equation}
\tilde{\nabla}_\theta J(\theta) = F(\theta)^{-1} \nabla_\theta J(\theta), \quad F(\theta) = \mathbb{E}_{s,a} \left[ \nabla_\theta \log \pi_\theta(a|s) \nabla_\theta \log \pi_\theta(a|s)^\top \right].
\end{equation}
Why does NPG recover policy iteration? Standard gradient ascent is sensitive to parameterization: it takes the steepest step in Euclidean parameter space, where units depend on how the policy is parameterized. NPG takes the steepest step in distribution space (measured by KL-divergence), which is invariant to reparameterization. In the tabular case, \citet[Theorem~2]{Kakade2001} proves that this geometric correction aligns the gradient exactly with the greedy policy $\tilde{\pi}$ from policy iteration. With step size 1 (and exact estimation), NPG performs one full Newton step; with smaller step sizes, it performs damped Newton updates. This explains its rapid convergence: NPG approximates the quadratic convergence of finding a fixed point rather than the linear convergence of hill-climbing.

The RL objective $J(\theta)$ is non-convex in $\theta$. For economists trained to distrust gradient methods on non-convex objectives, the natural concern is convergence to spurious local optima. For softmax policies, this concern is unfounded. The landscape is ``benign'' in a precise sense. \citet{agarwal2021theory} prove that $J(\theta)$ satisfies a \emph{gradient domination} condition (also called Polyak-\L{}ojasiewicz, or PL). The PL condition has the same functional form as the strong convexity condition for guaranteeing linear convergence of gradient descent, but it applies to non-convex functions: whenever $\|\nabla J(\theta)\|$ is small, the policy must be near-optimal. Formally, the sub-optimality $J(\pi^*) - J(\pi_\theta)$ is bounded by a constant times $\|\nabla J(\theta)\|^2$. The implication is immediate: any point where the gradient vanishes is globally optimal. The non-convex landscape has no false peaks, no spurious local maxima. Gradient ascent cannot get trapped.

\citet{mei2020} sharpen this result for softmax parameterization, proving explicit convergence rates.\footnote{However, ``no spurious local optima'' does not imply ``easy optimization.'' The landscape is dominated by vast plateaus (saddle points) where gradients vanish. Without sufficient exploration, the probability of visiting relevant states decays exponentially with the horizon, rendering the gradient exponentially small. Global convergence requires the starting distribution to have adequate coverage relative to the optimal policy's visitation distribution, formalized as the ``distribution mismatch coefficient'' by \citet{agarwal2021theory}. The Natural Policy Gradient addresses this by preconditioning with the Fisher Information Matrix, making the update direction covariant: invariant to invertible linear transformations of the parameter space. This standardizes units across parameters, preventing stalling on plateaus caused by poor parameter scaling. \citet{li2022softmax} make this quantitatively precise: vanilla softmax policy gradient requires iterations doubly exponential in the effective horizon $1/(1-\gamma)$ because score functions are exponentially small in directions corresponding to suboptimal actions.}

The Natural Policy Gradient achieves more: \emph{dimension-free} convergence. Standard gradient ascent on $J(\theta)$ has a convergence rate that depends on the smoothness constant, which scales with $|\mathcal{S}|$. NPG circumvents this by preconditioning with the Fisher information matrix $F(\theta)^{-1}$. The mechanism is that the state-visitation distribution $d^{\pi_\theta}(s)$ appears in both $\nabla J(\theta)$ and $F(\theta)$; when computing $F^{-1} \nabla J$, these terms cancel analytically. The resulting update rule is equivalent to soft policy iteration and converges at rate $O(1/(1-\gamma)^2 \epsilon)$, independent of $|\mathcal{S}|$ and $|\mathcal{A}|$ \citep{xiao2022convergence}.

\subsubsection{Trust Region Methods}

NPG requires computing and inverting the Fisher information matrix $F(\theta)$, which scales as $O(d^2)$ in parameters and is impractical for neural networks. TRPO \citep{Schulman2015} approximates the natural gradient using conjugate gradient methods\footnote{Conjugate gradient is an iterative method for solving linear systems $Ax = b$ without forming $A$ explicitly, requiring only matrix-vector products $Av$. With $k$ iterations it costs $O(kd)$ versus $O(d^3)$ for direct inversion, making it feasible for neural networks with millions of parameters.} without forming $F$ explicitly, and enforces trust regions via line search.\footnote{Trust region methods build on the performance difference lemma: for any two policies $\pi$ and $\pi'$, $J(\pi') - J(\pi) = \frac{1}{1-\gamma} \mathbb{E}_{s \sim d^{\pi'}} [ \sum_a \pi'(a|s) A^\pi(s,a) ]$, where $A^\pi(s,a) = Q^\pi(s,a) - V^\pi(s)$ is the advantage function. This identity bounds how much policy improvement is possible and motivates constraining updates to regions where advantage estimates remain accurate \citep{Kakade2002}.} \citet{shani2020} prove convergence for adaptive trust region methods that adjust the constraint radius dynamically. PPO \citep{Schulman2017} simplifies further by replacing the hard KL constraint with a clipped surrogate objective, trading theoretical guarantees for computational simplicity.

\subsection{Hybrid Methods}
\label{sec:actor_critic}

REINFORCE estimates policy gradients from sample returns (unbiased, high variance); TD methods use bootstrapped targets $r + \gamma V(s')$ (lower variance, biased when $V$ is approximate). Actor-critic methods combine both. The critic estimates the value function, the actor updates the policy using the critic's estimates.

\subsubsection{Actor-Critic Architecture and Two-Timescale Convergence}

The theoretical foundation is two-timescale stochastic approximation \citep{konda2000}, building on the two-timescale ODE convergence theory of \citet{borkar1997}.\footnote{The original analysis of \citet{konda2000} uses the average-cost formulation with TD error $\delta_t = c(X_t, U_t) - \Lambda + V(X_{t+1}) - V(X_t)$, where $\Lambda$ is the average cost. The discounted variant presented here follows by replacing the average-cost baseline with $\gamma V(s')$. A related but distinct ODE stability framework for single-timescale stochastic approximation, including Q-learning and TD learning, appears in \citet{borkar2000}.} Run two concurrent learning processes:
\begin{align}
\text{Critic (fast):} \quad & \theta_{t+1} = \theta_t + \alpha_t^{(c)} \delta_t \nabla_\theta \hat{V}(s_t; \theta_t), \\
\text{Actor (slow):} \quad & \omega_{t+1} = \omega_t + \alpha_t^{(a)} \delta_t \nabla_\omega \log \pi_\omega(a_t|s_t),
\end{align}
where $\delta_t = r_t + \gamma \hat{V}(s_{t+1}; \theta_t) - \hat{V}(s_t; \theta_t)$ is the TD error. The critic updates the value function parameters $\theta$; the actor updates the policy parameters $\omega$.

Convergence requires the critic to learn faster than the actor:
\begin{equation}
\lim_{t \to \infty} \frac{\alpha_t^{(a)}}{\alpha_t^{(c)}} = 0, \quad \text{with both satisfying Robbins-Monro conditions.}
\end{equation}
Under this separation, the actor sees a quasi-stationary critic: from the actor's perspective, the critic provides approximately correct value estimates at each step.\footnote{The two-timescale structure is analogous to nested optimization in structural estimation, where an inner loop solves for equilibrium given parameters and an outer loop searches over parameters. The critic's inner loop (policy evaluation) must converge before the actor's outer loop (policy improvement) takes a step. \citet{wu2020twotimescale} provide finite-time convergence rates ($\tilde{O}(\epsilon^{-2.5})$ sample complexity) for two-timescale actor-critic with linear approximation, and \citet{tian2023} establish analogous rates for single-timescale actor-critic with multi-layer neural networks.} The actor's updates are then approximately unbiased policy gradient steps. \citet{konda2000} prove convergence to a stationary point of $J(\omega)$ (i.e., $\nabla J(\omega) \to 0$).

Convergence requires a structural condition on the critic. The critic's feature vectors must span the actor's score functions $\nabla_\omega \log \pi_\omega(a|s)$, so that the critic's approximation error lies orthogonal to the policy gradient direction. Under this compatibility condition, the critic's projection error does not bias the actor's gradient estimates.\footnote{The compatible function approximation theorem first appears in \citet{SuttonMcAllester2000} and is the key structural requirement in \citet{konda2000}. It constrains the critic architecture to be ``compatible'' with the actor parameterization --- the same condition that makes the natural policy gradient equal to the critic's weight vector in \citet{Kakade2002}.}

A2C (Advantage Actor-Critic) is the synchronous variant: collect a batch of transitions, compute TD errors, and update both networks. A3C \citep{mnih2016a3c} parallelizes this across multiple workers updating a shared parameter server asynchronously.\footnote{Parallel workers decorrelate gradient estimates by exploring different parts of the state space simultaneously, removing the need for an experience replay buffer. Rigorous convergence theory for A3C's lock-free asynchronous parameter updates remains an open problem; existing analyses of asynchronous stochastic approximation \citep{qu2020async} address classical asynchrony (different state-action pairs updated at different times on a single trajectory), not the parallel-worker gradient setting.}

\subsubsection{Entropy Regularization and Soft Actor-Critic}

SAC (Soft Actor-Critic) \citep{Haarnoja2018} extends the actor-critic framework with entropy regularization, building on the soft Q-learning algorithm of \citet{Haarnoja2017}. The agent maximizes the entropy-augmented objective:
\begin{equation}
J_\tau(\theta) = \mathbb{E}_{\pi_\theta}\left[\sum_{t=0}^\infty \gamma^t \left(R_t + \tau \mathcal{H}(\pi_\theta(\cdot|s_t))\right)\right],
\end{equation}
where $\mathcal{H}(\pi) = -\sum_a \pi(a) \log \pi(a)$ is the entropy and $\tau > 0$ is the temperature parameter. \citet{geist2019regularized} later provided the unifying theoretical framework, showing that entropy regularization converts the Bellman optimality operator's non-smooth hard max into a smooth log-sum-exp, and that the resulting soft Bellman operator remains a $\gamma$-contraction, preserving the convergence guarantees of standard dynamic programming.\footnote{The optimal policy under entropy regularization is $\pi^*(a|s) \propto \exp(Q^*(s,a)/\tau)$, which is precisely the \citet{McFadden1974} logit choice probability with systematic utility $Q^*(s,a)$ and scale parameter $\tau$. The entropy-regularized value function is the log-sum-exp of Q-values, corresponding to the inclusive value (log-sum) operator in nested logit models. This equivalence between the soft-control framework and dynamic discrete choice models with EV1 taste shocks is developed formally in \citet{RustRawat2026}, Appendix~A.}

Entropy regularization also addresses the deadly triad directly. By maintaining policy stochasticity, the behavior policy used for data collection remains close to the target policy being optimized. This reduces the distribution mismatch between the stationary distribution under the behavior policy and the update targets, mitigating the off-policy instability leg of the triad. \citet{cen2022fast} formalize a second benefit: entropy regularization accelerates convergence of the natural policy gradient from $O(1/\epsilon)$ to $O(\log(1/\epsilon))$, providing a precise sense in which smoothing the policy landscape aids optimization. The actor-critic structure separates identification from optimization: the critic solves a regression problem (estimate $V^\pi$ from data), while the actor solves an optimization problem (improve $\pi$ using the estimated values).

\subsubsection{Error Amplification Under Approximate Value Functions}

Two questions remain important: how does approximation error propagate to policy quality, and how does computational complexity scale with problem size? \citet{singh1994} bound the policy degradation from value function errors (an independent derivation appears in \citealt{bertsekastsitsiklis1996}, Proposition~6.1). If $\hat{V}$ approximates $V^*$ with error $\|\hat{V} - V^*\|_\infty \leq \epsilon$, and $\hat{\pi}$ is the greedy policy with respect to $\hat{V}$, then:
\begin{equation}
\|V^* - V^{\hat{\pi}}\|_\infty \leq \frac{2\gamma}{1-\gamma} \epsilon.
\label{eq:singh_yee}
\end{equation}
At $\gamma = 0.99$, the amplification factor is $2 \cdot 0.99 / 0.01 = 198$. A 1\% error in value function approximation yields at most 198\% error in policy value.\footnote{The Singh-Yee bound is worst-case and not tight in general; \citet{farahmand2010} derive tighter bounds under smoothness assumptions on the MDP. The amplification factor $2\gamma/(1-\gamma)$ is a sensitivity analysis: it quantifies how errors in the ``inputs'' (value estimates) propagate to ``outputs'' (policy quality), analogous to errors-in-variables bias in regression. The discount factor $\gamma$ controls sensitivity; more patient agents face larger amplification.} This bound is pessimistic but finite: approximate value functions do not cause unbounded policy degradation.

\subsubsection{Sample Complexity of Planning}

Classical dynamic programming complexity scales with the state space size $|\mathcal{S}|$. For problems like Go, where $|\mathcal{S}| \approx 10^{170}$, exact computation is impossible. \citet{kearns2002} prove that with access to a generative model\footnote{A ``generative model'' in RL is a simulator that, given any state-action pair $(s,a)$, returns a sampled next state $s' \sim P(\cdot|s,a)$ and reward $r$. This is unrelated to ``generative models'' in machine learning (GANs, diffusion models) or ``generative processes'' in Bayesian econometrics. The distinction matters: planning with a generative model is strictly easier than learning from a single trajectory, because the agent can query arbitrary states rather than following a sequential path.} (a simulator that samples transitions from any state-action pair), near-optimal planning is possible with \emph{no dependence on $|\mathcal{S}|$}. The cost is exponential dependence on the effective horizon $H = \log(R_{\max}/(\epsilon(1-\gamma))) / \log(1/\gamma)$: the sparse sampling algorithm requires $O((|\mathcal{A}|/\epsilon)^{H})$ simulator calls.\footnote{The sparse sampling algorithm builds a random tree of depth $H$ from the current state, sampling $C$ successor states per action at each node, then estimates values by averaging leaf rewards back up the tree. Its running time is $O((C|\mathcal{A}|)^H)$, which is exponential in $H$ but entirely independent of $|\mathcal{S}|$. \citet{kearns2002} also prove a lower bound of $\Omega(2^H)$ generative model calls for any planning algorithm, so exponential horizon dependence is unavoidable in the worst case.} For $\gamma$ near 1, $H \approx (1-\gamma)^{-1} \log(R_{\max}/\epsilon)$, so the method is practical only for short effective horizons or moderate discount factors. The key insight is the tradeoff: classical DP scales linearly in $|\mathcal{S}|$ but polynomially in $H$; sparse sampling eliminates state-space dependence at the cost of exponential horizon dependence. This explains why MCTS succeeds in large state spaces with bounded lookahead.

The minimax-optimal sample complexity for planning with a generative model, when queries to arbitrary state-action pairs are permitted, is $\Theta(|\mathcal{S}||\mathcal{A}|/((1-\gamma)^3\epsilon^2))$ \citep{azar2013}. This bound scales linearly in $|\mathcal{S}|$ but polynomially in $1/(1-\gamma)$, the opposite regime from sparse sampling. \citet{agarwalKakadeYang2020} show that the plug-in model-based approach (learn $\hat{P}$ from samples, then plan with $\hat{P}$) achieves this minimax rate, establishing that model-based RL is statistically optimal. \citet{li2024breaking} further tighten this result by breaking the $|\mathcal{S}||\mathcal{A}|/(1-\gamma)^2$ sample-size barrier, showing that the minimax rate is achievable with total sample size as low as $|\mathcal{S}||\mathcal{A}|/(1-\gamma)$.\footnote{Recent extensions push these results beyond standard MDPs. \citet{clavier2024robust} study the robust MDP setting where the agent must plan under model uncertainty with sa-rectangular or s-rectangular uncertainty sets, establishing minimax rates for robust policy optimization. \citet{wang2025cvar} extend the analysis to risk-sensitive objectives under the iterated CVaR criterion.}

\subsection{Fundamental Tradeoffs}
\label{sec:tradeoffs}

The choice between methods involves distinct tradeoffs rooted in DP structure. Value-based methods target $Q^*$ via the Bellman contraction \citep{szepesvari2010}. In the tabular case convergence is guaranteed, but function approximation introduces the deadly triad. Policy-based methods optimize $\pi_\theta$ directly. Modern theory \citep{agarwal2021theory} establishes global convergence for softmax policies, with high variance as the practical weakness rather than local traps. Actor-critic methods combine both \citep{konda2000}, using the critic for low-variance value estimates while the actor inherits policy gradient's global convergence. Each family traces to DP foundations.

Four additional trade-offs pervade reinforcement learning. First, \emph{exploration versus exploitation}: should the agent act on its current best estimate or gather information to improve future decisions? \citet{LaiRobbins1985} establish the fundamental lower bound: any consistent policy must incur regret at least logarithmic in the number of periods. Naive exploration ($\varepsilon$-greedy) requires samples exponential in the horizon; strategic exploration (UCB, optimism in the face of uncertainty) reduces this to polynomial \citep{auer2002}, formalizing the value of targeted experimentation.\footnote{The exploration-exploitation tradeoff is the subject of the bandits chapter, where the multi-armed bandit framework provides the sharpest analysis. In full MDPs, exploration is harder because the agent must learn not just reward distributions but also transition dynamics, compounding the information requirement.}

Second, \emph{model-based versus model-free}: model-based methods learn a transition model $\hat{P}(s'|s,a)$ and plan with it \citep{sutton1990}; model-free methods learn value functions or policies directly from transitions. The Dyna architecture \citep{sutton1990} bridges these by generating simulated experience from the learned model to supplement real transitions. Model-based methods are sample-efficient (each transition updates the entire model, which improves value estimates for all states) but suffer asymptotic bias if the model class is misspecified; model-free methods are asymptotically unbiased but sample-inefficient, using each transition for a single gradient step.\footnote{Model misspecification in RL is the analog of omitted variable bias in econometrics: if the learned model omits relevant state variables or misspecifies the functional form of transitions, the resulting policy is biased regardless of sample size. \citet{moerland2023} provide a comprehensive survey of model-based RL, analyzing the model-bias versus sample-efficiency tradeoff across method families.}

Third, \emph{on-policy versus off-policy} \citep[Ch.~5--7]{sutton2018}: on-policy methods (SARSA, REINFORCE) learn from data generated by the current policy, ensuring stability but discarding past experience; off-policy methods (Q-learning, DQN) reuse stored experience via replay buffers, gaining sample efficiency but risking the instabilities of the deadly triad.

Fourth, \emph{bias versus variance in advantage estimation}: REINFORCE uses the full Monte Carlo return (unbiased but high variance); actor-critic methods \citep{konda2000} use bootstrapped TD targets (low variance but biased by the critic's approximation error). Generalized Advantage Estimation in PPO \citep{Schulman2015} interpolates between these extremes via a parameter $\lambda \in [0,1]$, where $\lambda = 1$ recovers Monte Carlo returns and $\lambda = 0$ recovers one-step TD. The value function baseline is the variance-minimizing choice, motivating the actor-critic architecture as a bias-variance compromise.

\subsection{Conclusion}
\label{sec:ch2_conclusion}

The central insight is that RL algorithms are not mysterious. They are asymptotic approximations to classical dynamic programming operators, justified by the mathematics of contractions, stochastic approximation, and gradient domination. Value iteration becomes Q-learning \citep{WatkinsDayan1992, tsitsiklis1994} when expectations are replaced by single samples. Policy iteration becomes the natural policy gradient \citep{Kakade2001, agarwal2021theory} when the greedy improvement step is approximated by gradient ascent, and NPG recovers PI exactly in the tabular case. The stochastic approximation framework, from the foundational work of \citet{Robbins1952} through the ODE method of \citet{borkar2000}, guarantees that under appropriate step-size conditions, noisy iterates converge to the same fixed points as their deterministic counterparts. Reinforcement learning is not a departure from dynamic programming but an extension of it. Tabular RL and RL with linear function approximation is well understood. However, deep RL is not as well understood. It seems to work for some cases but fails in others, and we do not have a firm handle on why. But in some sense, we do know that when it works it works due to the underlying Bellman principle. For these reasons convergence in RL remains and open issue and successes in RL remain brittle and case-specific.

